\documentclass[11pt]{article}

% ============================================================
% Distance-Constrained Dynamic Subtraction Games
% Overleaf-ready methodology / research record
% Compiler: pdfLaTeX
% ============================================================

\usepackage[margin=1in]{geometry}
\usepackage[T1]{fontenc}
\usepackage[utf8]{inputenc}
\usepackage{lmodern}
\usepackage{microtype}
\usepackage{amsmath,amssymb,amsthm,mathtools}
\usepackage{booktabs}
\usepackage{longtable}
\usepackage{array}
\usepackage{enumitem}
\usepackage{graphicx}
\usepackage{xcolor}
\usepackage{hyperref}
\usepackage{cleveref}

\hypersetup{
    colorlinks=true,
    linkcolor=blue,
    urlcolor=blue,
    citecolor=blue,
    pdftitle={Distance-Constrained Dynamic Subtraction Games},
    pdfsubject={Research Question, Methodology, and Computational Findings}
}

\setlist[itemize]{leftmargin=2em}
\setlist[enumerate]{leftmargin=2em}
\setlength{\parindent}{0pt}
\setlength{\parskip}{0.55em}
\allowdisplaybreaks

\newtheorem{theorem}{Theorem}
\newtheorem{conjecture}{Conjecture}
\newtheorem{proposition}{Proposition}
\newtheorem{lemma}{Lemma}
\newtheorem{observation}{Observation}

\newcommand{\game}{D(k,m)}
\newcommand{\mex}{\operatorname{mex}}
\newcommand{\Ppos}{\mathcal{P}}
\newcommand{\Npos}{\mathcal{N}}

\title{\textbf{Distance-Constrained Dynamic Subtraction Games}\\[0.4em]
\large Research Question, Methodology, and Computational Findings to Date}
\author{Research Team}
\date{\today}

\begin{document}

\maketitle

\begin{center}
\begin{tabular}{@{}ll@{}}
\textbf{Project area:} & Combinatorial Game Theory / Discrete Mathematics \\
\textbf{Current main case:} & $D(k,1)$ \\
\textbf{Status:} & Computational discovery and validation completed for the main observed regimes; \\
& proof-oriented full-state analysis is beginning.
\end{tabular}
\end{center}

\tableofcontents
\newpage

% ============================================================
\section{Research Question}
% ============================================================

\subsection{Primary Research Question}

\begin{quote}
\textbf{How does imposing a minimum separation between consecutive move sizes change the periodic structure of losing positions in a dynamic subtraction game?}
\end{quote}

More specifically, for the game $D(k,m)$, we investigate:

\begin{quote}
For fixed move bound $k$ and move-separation parameter $m$, what are the P-positions (previous-player winning / next-player losing positions), do they become eventually periodic, and can their periods and residue classes be classified as functions of $k$ and $m$?
\end{quote}

\subsection{Current Focused Question}

The main investigation is now concentrated on $m=1$:

\begin{quote}
\textbf{For $D(k,1)$, can the observed affine families of periods and P-position residues be explained and proved from the full state structure $(n,s)$?}
\end{quote}

A stronger long-term question is whether there is a parameterized family of regimes of the form
\begin{equation}
    p(k)=2k+A+2,
    \qquad
    P_k\equiv\{0,A,k+A+1\}\pmod{p(k)},
\end{equation}
with mathematically describable values of $A$ and corresponding ranges of $k$.

% ============================================================
\section{Definition of the Game}
% ============================================================

We define a family of impartial combinatorial games $D(k,m)$.

\begin{itemize}
    \item A single heap initially contains $n$ counters.
    \item On each turn, a player removes an integer number
    \[
        r\in\{1,2,\ldots,k\},
    \]
    subject to $r\le n$.
    \item If the previous move removed $s$ counters, the next move must satisfy
    \[
        |r-s|>m.
    \]
    \item The first move is unrestricted by a previous move.
    \item Normal play is used: a player with no legal move loses, even if counters remain in the heap.
\end{itemize}

The parameter $k$ controls the largest possible removal. The parameter $m$ controls how different the next move must be from the previous move.

Important cases studied so far are:
\begin{itemize}
    \item $m=0$: consecutive moves cannot be equal; used as a known/control baseline.
    \item $m=1$: consecutive move sizes must differ by at least $2$; this is the main research case.
    \item $m=2$: exploratory comparison/generalization.
\end{itemize}

% ============================================================
\section{State-Space Model}
% ============================================================

Because legal moves depend on the previous move, the heap size $n$ alone is not enough to determine the game state.

After the first move, a state is represented as
\begin{equation}
    (n,s),
\end{equation}
where $n$ is the number of counters remaining and $s$ is the number of counters removed on the previous move.

From $(n,s)$, a move of size $r$ is legal when
\begin{equation}
    1\le r\le\min(k,n)
\end{equation}
and
\begin{equation}
    |r-s|>m.
\end{equation}
The resulting state is
\begin{equation}
    (n-r,r).
\end{equation}

This enlarged state representation converts the history-dependent rule into a memoryless directed game graph.

% ============================================================
\section{Sprague--Grundy Recurrence}
% ============================================================

For every state $(n,s)$, define the Sprague--Grundy value $g(n,s)$. The legal child states are $(n-r,r)$ for every legal $r$. Therefore,
\begin{equation}
\boxed{
    g(n,s)=
    \mex\left\{
        g(n-r,r):
        1\le r\le\min(k,n),\ |r-s|>m
    \right\}.
}
\end{equation}

Here, $\mex$ is the minimum excluded nonnegative integer. A state is a P-position exactly when
\begin{equation}
    g(n,s)=0.
\end{equation}

For the unrestricted starting position, define
\begin{equation}
\boxed{
    g_0(n)=
    \mex\left\{
        g(n-r,r):1\le r\le\min(k,n)
    \right\}.
}
\end{equation}
Thus an initial heap size $n$ is a starting P-position exactly when
\begin{equation}
    g_0(n)=0.
\end{equation}

The binary starting P/N sequence is recorded as
\begin{equation}
    b_n=
    \begin{cases}
        1, & g_0(n)=0,\\
        0, & g_0(n)\ne0.
    \end{cases}
\end{equation}

% ============================================================
\section{Reference Solver}
% ============================================================

A direct reference solver was implemented first. For each $n$, each previous move $s$, and each possible removal $r$, the solver explicitly tests whether
\begin{equation}
    |r-s|>m.
\end{equation}
It collects the SG values of every legal child and computes the mex directly.

This solver is intentionally simple. Its main purpose is not speed but independent correctness checking.

% ============================================================
\section{Optimized Solver}
% ============================================================

The optimized solver uses the fact that, for fixed $n$, the potential child value associated with removing $r$ is
\begin{equation}
    v_r=g(n-r,r).
\end{equation}

For an unrestricted starting state, all valid $r$ are available. For a state with previous move $s$, only the interval
\begin{equation}
    s-m\le r\le s+m
\end{equation}
is forbidden, after clipping to the legal move range.

The optimized method therefore:
\begin{enumerate}
    \item computes the SG values $v_r$ for all possible removals;
    \item counts the occurrences of each child SG value;
    \item computes the unrestricted mex;
    \item for each previous move $s$, temporarily removes the at-most-$(2m+1)$ forbidden move contributions;
    \item updates the mex after those removals;
    \item stores $g(n,s)$; and
    \item restores the removed contributions before processing the next $s$.
\end{enumerate}

For fixed small $m$, the approximate time complexity is
\begin{equation}
    O\!\left(Nk(m+1)\right).
\end{equation}
For the main case $m=1$, this is approximately
\begin{equation}
    O(Nk)
\end{equation}
up to constant factors. The full stored SG table uses approximately
\begin{equation}
    O(Nk)
\end{equation}
memory.

% ============================================================
\section{Independent Verification}
% ============================================================

The optimized implementation was not accepted solely because it produced plausible sequences. A separate slow reference implementation was used for automated self-tests on small parameter values.

The two implementations were compared across small values of $k$, $m$, and $n$, checking both the full state values $g(n,s)$ and the unrestricted starting values $g_0(n)$.

The C++ research implementation passed these self-tests before deep scans were performed. This verification step was included to reduce the possibility that an optimization error could create false mathematical patterns.

% ============================================================
\section{Python and C++ Research Workflow}
% ============================================================

The project began with a Python research implementation because Python allowed rapid development and inspection of the recurrence, P-position sequences, periods, and residues.

After the computational range became large, the same research engine was ported to C++17 for faster deep scans. The C++ implementation retained:
\begin{itemize}
    \item SG computation;
    \item starting P/N sequence generation;
    \item candidate eventual-period detection;
    \item residue extraction;
    \item affine period-law mining;
    \item contiguous affine-regime detection;
    \item unusual-case detection;
    \item optional full-state cycle detection;
    \item CSV output;
    \item automated research reports;
    \item independent self-tests.
\end{itemize}

Python remains useful for the present proof-discovery stage because it allows flexible analysis of the full $(n,s)$ state structure and symbolic pattern exploration.

% ============================================================
\section{Period Detection Method}
% ============================================================

For each fixed pair $(k,m)$, the starting P/N sequence
\begin{equation}
    b_0,b_1,\ldots,b_N
\end{equation}
was generated.

A period detector searched for candidate periods $p$ such that a sufficiently long computed tail satisfies
\begin{equation}
    b_n=b_{n+p}.
\end{equation}

The detector requires repeated copies of the candidate period and then searches backward to estimate the candidate preperiod.

When a candidate period is detected, the program records:
\begin{itemize}
    \item candidate preperiod;
    \item candidate period $p$;
    \item P-position density within the period;
    \item P-position residues modulo $p$;
    \item maximum observed SG values;
    \item a fingerprint of the repeating word.
\end{itemize}

\subsection{Important Limitation}

A finite repeating tail is \textbf{not} a mathematical proof of eventual periodicity. All period statements produced by this detector are treated as computational observations or conjectures until proved from the recurrence or certified by an appropriate exact-state argument.

% ============================================================
\section{Control and Comparative Experiments}
% ============================================================

The study compared $m=0$, $m=1$, and $m=2$.

\subsection{$m=0$}
For $m=0$, the rule becomes
\begin{equation}
    r\ne s.
\end{equation}
This repetition-free case provides a control/baseline. The program rediscovered highly structured period behavior in this case, supporting the validity of the computational approach while also demonstrating that known/control behavior should not be treated as a new result.

\subsection{$m=1$}
The $m=1$ case produced a mixture of simple periodic blocks, long periods, abrupt transitions, and especially clean affine regimes. It was therefore selected as the main research focus.

\subsection{$m=2$}
The $m=2$ case was explored as a generalization and comparison. It also showed complex behavior, but the current proof effort is intentionally concentrated on $m=1$ to keep the mathematical scope manageable.

% ============================================================
\section{Pattern-Mining Procedure}
% ============================================================

Two main automated pattern searches were used.

\subsection{Affine Laws Across Parameter Classes}
The program searched for relationships of the form
\begin{equation}
    p(k)=ak+b
\end{equation}
among sets of parameter values.

\subsection{Contiguous Affine Regimes}
A second detector was added to identify consecutive intervals of $k$ for which the same affine period formula holds. This was necessary because a search based only on global congruence classes can miss a formula that is valid only on one consecutive range.

The contiguous search led to the main observed $m=1$ regimes.

% ============================================================
\section{Main Computational Findings So Far}
% ============================================================

All statements in this section are \textbf{observed computational patterns unless explicitly stated otherwise}.

\subsection{First Affine Regime}
For
\begin{equation}
    42\le k\le52,
\end{equation}
the deep computations support
\begin{equation}
\boxed{p(k)=2k+56}
\end{equation}
with starting P-position residues
\begin{equation}
\boxed{
    P_k\equiv\{0,54,k+55\}\pmod{2k+56}.
}
\end{equation}
The P-position density within the candidate period is therefore
\begin{equation}
    \frac{3}{2k+56}.
\end{equation}

\subsection{Second Affine Regime}
For
\begin{equation}
    58\le k\le78,
\end{equation}
the deep computations support
\begin{equation}
\boxed{p(k)=2k+88}
\end{equation}
with starting P-position residues
\begin{equation}
\boxed{
    P_k\equiv\{0,86,k+87\}\pmod{2k+88}.
}
\end{equation}
The P-position density is
\begin{equation}
    \frac{3}{2k+88}.
\end{equation}

This regime contains 21 consecutive values of $k$, making it the current primary theorem candidate.

\subsection{Common Template}
The two regimes share the form
\begin{equation}
\boxed{p=2k+A+2}
\end{equation}
and
\begin{equation}
\boxed{
    P_k\equiv\{0,A,k+A+1\}\pmod p.
}
\end{equation}
The observed constants are
\begin{equation}
    A=54
\end{equation}
and
\begin{equation}
    A=86.
\end{equation}
Their difference is
\begin{equation}
    86-54=32.
\end{equation}
This numerical relationship is currently only suggestive. No recurrence or binary/2-adic explanation has been proved.

% ============================================================
\section{Observed Boundary Behavior}
% ============================================================

The clean affine laws do not continue universally. Immediately after the second regime, $k=79$ has an observed candidate period of
\begin{equation}
    166
\end{equation}
with a nonzero preperiod, while $k=80$ has a much larger observed candidate period
\begin{equation}
    947.
\end{equation}

This abrupt change suggests that the affine formulas belong to genuine parameter regimes whose boundaries must be explained.

Other observed simple blocks include a region around
\begin{equation}
    93\le k\le99
\end{equation}
with
\begin{equation}
    p=k+1
\end{equation}
and starting residue $\{0\}$.

At $k=100$, the observed period remains $101$, but a nonzero preperiod appears, showing that period length alone does not capture the full structure.

A further simple block was observed beginning at approximately
\begin{equation}
    117\le k\le120,
\end{equation}
again with
\begin{equation}
    p=k+1
\end{equation}
and residue $\{0\}$. The scan ended at $k=120$, so the upper boundary of this block has not yet been determined.

% ============================================================
\section{Deep Validation Experiment}
% ============================================================

The principal deep validation scan used
\begin{equation}
    m=1,
    \qquad
    35\le k\le120,
    \qquad
    N=500{,}000,
\end{equation}
with candidate periods searched up to
\begin{equation}
    150{,}000.
\end{equation}

The execution was divided into two independent ranges because of a runtime limit:
\begin{equation}
    35\le k\le80
\end{equation}
and
\begin{equation}
    81\le k\le120.
\end{equation}

Each $k$ is computed independently, so splitting the range does not change the recurrence or the resulting data.

The two principal affine regimes continued to appear through this much larger computational horizon. This provides substantially stronger computational evidence than the earlier shorter scans, while still not replacing proof.

% ============================================================
\section{False-Tail Detection and Stress Testing}
% ============================================================

One earlier shorter computation for $k=80$ suggested a candidate preperiod near
\begin{equation}
    99{,}685.
\end{equation}

After increasing the computational horizon to $N=500{,}000$, the candidate changed to
\begin{equation}
\boxed{
    \text{preperiod}=0,
    \qquad
    p=947.
}
\end{equation}

The observed residues were
\begin{equation}
    \{0,86,168,249,785,866\}\pmod{947}.
\end{equation}

This is an important methodological example: a finite sequence can contain a long tail that appears periodic without representing the true eventual behavior.

Therefore the project distinguishes between:
\begin{enumerate}
    \item finite computational observations;
    \item conjectured eventual behavior;
    \item mathematically proved behavior.
\end{enumerate}

The $k=80$ correction is retained as evidence for why deeper stress testing is necessary.

% ============================================================
\section{Other Notable Deep-Scan Observations}
% ============================================================

The $N=500{,}000$ scan also produced several large or irregular candidate periods. Examples include
\begin{equation}
    p(80)=947,
\end{equation}
\begin{equation}
    p(81)=597,
\end{equation}
and
\begin{equation}
    p(107)=1370.
\end{equation}

For $k=107$,
\begin{equation}
    \frac{p}{k}=\frac{1370}{107}\approx12.80.
\end{equation}

These cases demonstrate the complexity of the game, but they are not currently the primary proof targets. Clean parameterized regimes are considered more useful for theorem discovery than isolated large periods.

% ============================================================
\section{Interpretation of the Evidence}
% ============================================================

The current evidence suggests that $D(k,1)$ does not follow one universal simple period formula. Instead, $k$ appears to pass through distinct structural regimes.

Within some regimes:
\begin{itemize}
    \item the period is affine in $k$;
    \item the number of starting P-residues is fixed;
    \item the residues themselves are affine functions of $k$.
\end{itemize}

At regime boundaries, the behavior can change abruptly.

This motivates the working hypothesis:

\begin{conjecture}[Working hypothesis]
For $m=1$, the parameter $k$ is partitioned into structural regimes in which the period and starting P-position residues admit simple affine descriptions.
\end{conjecture}

This is a hypothesis, not yet a theorem.

% ============================================================
\section{Current Primary Theorem Candidate}
% ============================================================

The most attractive current finite-range theorem candidate is:

\begin{conjecture}[Primary finite-range candidate]
For every integer $k$ with
\[
    58\le k\le78,
\]
the starting P-positions of $D(k,1)$ are exactly the integers satisfying
\[
    n\equiv0,\ 86,\ k+87\pmod{2k+88}.
\]
\end{conjecture}

The computational evidence through $N=500{,}000$ strongly supports this statement. However, a symbolic proof cannot be obtained from the starting sequence alone. The proof is expected to require classification of the complete state $(n,s)$.

% ============================================================
\section{Next Phase: Full-State Structural Analysis}
% ============================================================

The next phase changes the purpose of computation. Instead of scanning larger $k$, Python will be used to examine
\begin{equation}
    g(n,s)=0
\end{equation}
throughout the complete state space.

For each selected $k$, define
\begin{equation}
    Z_{k,s}=\{n\bmod p:g(n,s)=0\},
\end{equation}
where $p$ is the observed starting period.

The analysis will:
\begin{enumerate}
    \item compute the full SG table;
    \item extract every zero-SG state $(n,s)$;
    \item reduce $n$ modulo the candidate period;
    \item compare the zero-state residue sets across $s$;
    \item search for affine boundaries and formulas involving $s$ and $k$;
    \item compare the same normalized structure across several $k$-values in the $58\le k\le78$ regime;
    \item formulate a candidate complete description of the P-states.
\end{enumerate}

The most useful outcome would be a compact description of a set
\begin{equation}
    \Ppos_k=\{(n,s):\text{specified modular and inequality conditions}\}.
\end{equation}

% ============================================================
\section{Planned Proof Strategy}
% ============================================================

Once a candidate full-state description $\Ppos_k$ is identified, the intended proof method is the standard P/N characterization.

\subsection{Condition 1: No P-to-P Move}
Show that if
\begin{equation}
    (n,s)\in\Ppos_k,
\end{equation}
then no legal move $r$ satisfying
\begin{equation}
    |r-s|>1
\end{equation}
produces another state in $\Ppos_k$.

\subsection{Condition 2: Every N-State Reaches P}
Show that every state outside $\Ppos_k$ has at least one legal removal $r$ such that
\begin{equation}
    (n-r,r)\in\Ppos_k.
\end{equation}

If both conditions hold, $\Ppos_k$ is exactly the set of P-positions. The unrestricted starting-position theorem can then be derived from the legal first moves into these classified states.

% ============================================================
\section{Stronger Long-Term Target}
% ============================================================

If the two observed affine regimes arise from the same mechanism, the stronger target is a family indexed by a parameter $t$.

The desired type of result would be:
\begin{equation}
    L_t\le k\le U_t
\end{equation}
implies
\begin{equation}
    p_t(k)=2k+A_t+2
\end{equation}
and
\begin{equation}
    P_k\equiv\{0,A_t,k+A_t+1\}\pmod{p_t(k)}.
\end{equation}

The main unanswered questions are:
\begin{itemize}
    \item What determines $A_t$?
    \item What determines the regime boundaries $L_t,U_t$?
    \item Why do the observed values $54$ and $86$ occur?
    \item Is there an infinite sequence of such regimes?
    \item Why do simple $p=k+1$ blocks reappear between more complex regimes?
\end{itemize}

These questions will only be promoted to formal conjectures when additional structural evidence is available.

% ============================================================
\section{Reproducibility and Research Records}
% ============================================================

For each computational experiment, the following should be retained:
\begin{itemize}
    \item source-code version;
    \item programming language and compiler/interpreter version;
    \item exact command or parameters;
    \item values of $m$, $k$, and $N$;
    \item maximum searched period;
    \item raw CSV output;
    \item generated report;
    \item date of run;
    \item whether the result is exploratory, deeply validated, or proved;
    \item any later run that contradicts or modifies an earlier interpretation.
\end{itemize}

The slow reference solver and automated self-tests should remain part of the codebase.

% ============================================================
\section{Current Status Summary}
% ============================================================

\subsection{Established Computationally Within Tested Ranges}
\begin{itemize}
    \item The recurrence and solver have been independently cross-checked on small cases.
    \item $m=0,1,2$ have been compared computationally.
    \item $m=1$ displays distinct structural regimes.
    \item The $42\le k\le52$ affine pattern survived deep testing to $N=500{,}000$.
    \item The $58\le k\le78$ affine pattern survived deep testing to $N=500{,}000$.
    \item Deep testing corrected an earlier misleading finite-tail interpretation at $k=80$.
\end{itemize}

\subsection{Not Yet Proved Symbolically}
\begin{itemize}
    \item Eventual periodicity for the observed $m=1$ regimes in a single parameterized symbolic proof.
    \item A general rule generating the constants $54,86,\ldots$.
    \item A formula for the boundaries of the affine regimes.
    \item An infinite family of affine regimes.
\end{itemize}

\subsection{Immediate Research Objective}
\begin{quote}
\textbf{Determine the complete zero-SG state structure $(n,s)$ for the $58\le k\le78$ regime and use it to formulate a proof-ready classification theorem.}
\end{quote}

% ============================================================
\section{Research Principle Going Forward}
% ============================================================

The computational program is being used to
\begin{equation}
\boxed{
\text{discover patterns}
\longrightarrow
\text{stress-test them}
\longrightarrow
\text{find counterexamples}
\longrightarrow
\text{formulate conjectures}.
}
\end{equation}

The next stage is to convert those conjectures into mathematical statements that can be proved directly from the rules of the game.

Large computations are therefore not the final result; they are the experimental tool used to discover the underlying combinatorial structure.

% ============================================================
\section{First Result from the Full-State Python Phase}
% ============================================================

After the methodology above was established, the next proof-oriented computation was performed in Python.

For $m=1$, define the \textbf{winning-move set}
\begin{equation}
    A_n=\{r\in\{1,\ldots,\min(k,n)\}:(n-r,r)\text{ is a P-state}\}.
\end{equation}

This gives an exact Boolean reformulation of the game.

The unrestricted starting position at heap size $n$ is P exactly when
\begin{equation}
\boxed{A_n=\varnothing.}
\end{equation}

For a state with previous move $s$, the only forbidden next removals are
\begin{equation}
    \{s-1,s,s+1\}
\end{equation}
after clipping to $\{1,\ldots,k\}$. Therefore,
\begin{equation}
\boxed{
    (n,s)\text{ is P}
    \iff
    A_n\subseteq\{s-1,s,s+1\}.
}
\end{equation}

This identity follows directly from the P/N recurrence: a state is P precisely when every move that would reach a P-state is forbidden.

This reformulation is useful because the proof search no longer needs arbitrary SG values. It can work with Boolean P/N rows and the sets $A_n$.

% ============================================================
\section{Independent Boolean Solver}
% ============================================================

A second Python solver was implemented using only P/N information.

For each $n$:
\begin{enumerate}
    \item it forms $A_n$ by checking which removals $r$ lead to a previously computed P-state $(n-r,r)$;
    \item it declares the unrestricted start at $n$ P if and only if $A_n$ is empty;
    \item for each previous move $s$, it declares $(n,s)$ P if and only if every member of $A_n$ lies inside the forbidden interval $[s-1,s+1]$.
\end{enumerate}

This Boolean solver does not calculate mex or nonzero SG values.

Its complete P/N output was compared against the original SG solver for the tested $58\le k\le78$ cases, and the two methods agreed on both:
\begin{itemize}
    \item unrestricted starting P/N positions;
    \item every full state $(n,s)$ checked.
\end{itemize}

This provides a second implementation of the part of the recurrence relevant to P-position classification.

% ============================================================
\section{Exact Finite-State Period Certificate}
% ============================================================

The Boolean recurrence has finite memory.

To compute row $n$, the game only needs states of the form
\begin{equation}
    (n-r,r),
    \qquad
    1\le r\le k.
\end{equation}
Therefore row $n$ depends only on the preceding $k$ rows.

Let $R_n$ denote the complete Boolean P/N row
\begin{equation}
    R_n=\bigl(P(n,1),P(n,2),\ldots,P(n,k)\bigr).
\end{equation}

Suppose that for some offset $t$ and candidate period $p$,
\begin{equation}
    R_{t+i}=R_{t+p+i}
\end{equation}
for every
\begin{equation}
    0\le i<k.
\end{equation}

Then the next rows must also agree:
\begin{equation}
    R_{t+k}=R_{t+p+k},
\end{equation}
because they are computed from identical preceding $k$-row states.

Repeating this argument inductively proves
\begin{equation}
\boxed{
    R_n=R_{n+p}\quad\text{for every }n\ge t.
}
\end{equation}

Thus equality of two $k$-row state windows is an \textbf{exact finite certificate of future P/N periodicity}, rather than a finite-tail heuristic.

% ============================================================
\section{Certificate Results for the $58\le k\le78$ Regime}
% ============================================================

The Python Boolean solver was used to test the candidate period
\begin{equation}
    p=2k+88
\end{equation}
for every integer
\begin{equation}
    58\le k\le78.
\end{equation}

For every one of the 21 parameter values, an exact repeated $k$-row P/N state window was found.

For
\begin{equation}
    65\le k\le78,
\end{equation}
the full Boolean state pattern repeats with period $2k+88$ beginning at
\begin{equation}
    t=0.
\end{equation}

For
\begin{equation}
    58\le k\le64,
\end{equation}
the full Boolean state pattern has a short transient, with the certificate beginning near $t=58$ or $59$.

\begin{table}[htbp]
\centering
\caption{Exact finite-state certificate data for $58\le k\le78$.}
\label{tab:certificates}
\small
\begin{tabular}{@{}rrrrl@{}}
\toprule
$k$ & $p=2k+88$ & Full-state cycle start $t$ & Prefix checked through & Starting P-residues \\
\midrule
58 & 204 & 58 & 115 & $\{0,86,145\}$ \\
59 & 206 & 59 & 117 & $\{0,86,146\}$ \\
60 & 208 & 59 & 118 & $\{0,86,147\}$ \\
61 & 210 & 59 & 119 & $\{0,86,148\}$ \\
62 & 212 & 59 & 120 & $\{0,86,149\}$ \\
63 & 214 & 58 & 120 & $\{0,86,150\}$ \\
64 & 216 & 59 & 122 & $\{0,86,151\}$ \\
65 & 218 & 0 & 64 & $\{0,86,152\}$ \\
66 & 220 & 0 & 65 & $\{0,86,153\}$ \\
67 & 222 & 0 & 66 & $\{0,86,154\}$ \\
68 & 224 & 0 & 67 & $\{0,86,155\}$ \\
69 & 226 & 0 & 68 & $\{0,86,156\}$ \\
70 & 228 & 0 & 69 & $\{0,86,157\}$ \\
71 & 230 & 0 & 70 & $\{0,86,158\}$ \\
72 & 232 & 0 & 71 & $\{0,86,159\}$ \\
73 & 234 & 0 & 72 & $\{0,86,160\}$ \\
74 & 236 & 0 & 73 & $\{0,86,161\}$ \\
75 & 238 & 0 & 74 & $\{0,86,162\}$ \\
76 & 240 & 0 & 75 & $\{0,86,163\}$ \\
77 & 242 & 0 & 76 & $\{0,86,164\}$ \\
78 & 244 & 0 & 77 & $\{0,86,165\}$ \\
\bottomrule
\end{tabular}
\end{table}

Despite the short full-state transient for $58\le k\le64$, the unrestricted starting P-position sequence was separately checked through the finite prefix needed to connect to the certified cycle.

For every
\begin{equation}
    58\le k\le78,
\end{equation}
the starting sequence is thereby exactly certified to have period
\begin{equation}
\boxed{2k+88}
\end{equation}
from $n=0$, with residues
\begin{equation}
\boxed{
    \{0,86,k+87\}.
}
\end{equation}

This is stronger than observing the same pattern through $N=500{,}000$: for each of these 21 fixed values of $k$, the finite-state recurrence supplies a certificate that forces the P/N pattern to continue indefinitely.

This remains a \textbf{computer-assisted finite-range certification}, not yet a symbolic human proof explaining simultaneously why every $k$ in the interval has this form or why the interval boundaries are $58$ and $78$.

% ============================================================
\section{Updated Immediate Objective}
% ============================================================

The project now has two complementary goals.

\subsection{Exact Finite-Range Result}
Preserve and independently verify the Boolean cycle certificates for all
\begin{equation}
    58\le k\le78.
\end{equation}
A small standalone verifier should eventually be produced so that the certificate logic can be inspected independently of the exploratory code.

\subsection{Symbolic Mathematical Result}
Use the winning-move sets $A_n$ to explain why
\begin{equation}
    p=2k+88
\end{equation}
and
\begin{equation}
    P_k\equiv\{0,86,k+87\}\pmod{2k+88}
\end{equation}
arise, and why the same structure holds exactly over the observed parameter regime.

The next pattern-mining target is therefore not the SG values themselves, but the structure of the sets $A_n$, especially the residues for which $A_n$ is empty or is contained in three consecutive move sizes.

% ============================================================
\section{Conclusion and Immediate Next Step}
% ============================================================

The research has progressed from exploratory computation to an exact finite-state certification for the principal $58\le k\le78$ regime. The most important next step is to convert the visible full-state geometry into a compact symbolic classification.

In particular, the next Python analysis should attempt to express the zero-SG state diagram as a small collection of affine families, for example conditions of the form
\begin{equation}
    n=s+C,
\end{equation}
\begin{equation}
    n=2s+C,
\end{equation}
or bounded affine regions whose endpoints depend on $k$.

A successful classification would provide the structure needed to prove the two P/N conditions directly and explain the observed formula
\begin{equation}
    p=2k+88,
    \qquad
    P_k\equiv\{0,86,k+87\}\pmod{2k+88}.
\end{equation}

\end{document}
